\section{An analytic witness for the overlap correction}\label{app:witness}

We retain the local sign witness of \cite[Section 5.3]{phase} to show that \(\gamma\) in Theorem~\ref{thm:overlap} is not identically zero. This calculation analytically differentiates the native map. It is not a numerical native-step evaluation, sampled trajectory, parameter scan or new continuation problem.

Take the canonical family \(z(t)=(1,e^{it},e^{it})\), \(|t|\le D_0\). Its amplitudes are one and its phase width is \(|t|\), so every input is admitted. The intensities and hence \(M\) are constant on this family. At \(t=0\), its prestage is the positive vector
\begin{equation}
w_i=1+\eps(k_i-1),\qquad
w=\left(\frac{9999}{10000},1,\frac{10001}{10000}\right),\qquad \sum_i w_i=3.
\label{eq:witnessw}
\end{equation}
Differentiating the prestage at zero gives \(i(2g,w_2-g,w_3-g)\), so its phase derivatives are
\begin{equation}
\tau_1=\frac{2g}{w_1},\qquad \tau_2=1-\frac g{w_2},\qquad
\tau_3=1-\frac g{w_3}.
\label{eq:witnesstau}
\end{equation}
Write \(d=1-9\lambda>0\), \(T=\sum_i\tau_i\). Differentiating the simultaneous kicks at the aligned prestage gives output phase derivatives
\begin{equation}
\sigma_i=\tau_i+3\lambda\sum_{j\ne i}(\tau_j-\tau_i)
=d\tau_i+3\lambda T.
\label{eq:witnesssigma}
\end{equation}
The output modulus derivatives vanish: the baseline vector is real positive and its derivative is imaginary. Removing the output channel-1 phase therefore gives derivatives \(i w_i(\sigma_i-\sigma_1)\) for the next canonical representative. The current conjugate input derivatives are \((0,-i,-i)\). Thus
\begin{equation}
B_{\rm sec}(0)=3,\qquad
\im B_{\rm sec}'(0)=\sum_i w_i\sigma_i-3\sigma_1-(w_2+w_3).
\label{eq:witnessB}
\end{equation}
Since \(\sum_iw_i\tau_i=w_2+w_3\), substitution cancels the \(T\) terms and yields
\begin{equation}
\im B_{\rm sec}'(0)=-9\lambda(w_2+w_3)-3d\tau_1,
\qquad
\boxed{\gamma'(0)=-3\lambda(w_2+w_3)-d\tau_1<0.}
\label{eq:witnesssign}
\end{equation}
The established right-half-plane branch justifies differentiating its argument and dividing by \(B_{\rm sec}(0)=3\). Therefore \(\gamma(t)\ne0\) for sufficiently small positive \(t\), and \(\alpha_{\rm ov}\ne\vartheta\) there. The conclusion is local nonidentity, not a classification of all states where the angles agree or a verdict on endpoint-difference representations.
